\section{Boundary, Wall and Timestep Treatment}

The production boundary policy includes a Regional2D-forced inlet, open outlet, lateral open-ocean boundaries, atmospheric top, terrain wall, barrier wall and isotropic damping zones. The reflection benchmark distinguishes reflective control behavior from production outlet and lateral behavior:

\begin{itemize}
  \item reflective control: \(K_r=0.8199788627\), \(RE=0.672400004\);
  \item production outlet: \(K_r=0.0799979378\), \(RE=0.00640000004\);
  \item production lateral: \(K_r=0.1000003891\), \(RE=0.00999999974\).
\end{itemize}

The wall-function baseline uses continuous Spalding treatment for turbulent viscosity, while \texttt{kqRWallFunction} and \texttt{omegaWallFunction} are retained for turbulence quantities. Runtime \(y^+\) evidence is accepted for the G6 baseline. Formal wall-function/mesh convergence remains post-G6.

The timestep policy is accepted for the G6 baseline: repository configuration defines pre-run caps and minimum timestep acceptance, while OpenFOAM controls internal Courant and interface-Courant adaptation. Formal timestep convergence and an optional external rollback/retry supervisor remain post-G6, preserving issue \#62.
